\section{Discussion}
In this project, we implemented a strongly modular framework for basic gymnasium environments, which can be easily extended to other environments like Hockey. This enabled us to provide a deep evaluation based on multiple environments and for multiple hyperparameter configurations. Furthermore, the framework allows us to implement and compare multiple agents with each other based on the implemented tournament (\autoref{sec:final-tournament}). The presented evaluation results, show how big the influence of better buffers, rewards, exploration noise on different model types and understanding the game can be.

Firstly, we see that prioritization or balancing of buffers works well on well-known benchmarks, but cannot be applied to the Hockey environment. This is strongly related to the sparse reward and big buffers required for the environment. Instead, the exploration with pink noise is a good choice, because we observe the highest final reward and fastest convergence in our evaluation.

Furthermore, \ac{SAC} is shown to converge almost in half of the steps compared to our best \ac{TD3}. This effect is enforced by the alpha-tuning, which dynamically weights the exploration noise. Together with our custom reward containing penalties for aggressive behavior and refusal to play, we can provide almost equal performing \ac{SAC} and \ac{TD3} implementations.

Moreover, our hybrid model is highly adaptable and based on the previously trained agents, with the aim of compensating for the weaknesses of the other model. Thus, it can use the more aggressive \ac{TD3} agent for shooting and the passive \ac{SAC} for defending the goal. However, we saw in our training and evaluation experiments that it tends to overfit fast to the provided opponents. In addition, the graphical evaluation shows that \ac{SAC} and \ac{TD3} learn different policies. \ac{SAC} tends to shoot directly to the goal and act more passive. Instead, focuses \ac{TD3} on shooting over the horizontal walls by being more aggressive in the center of the field. Besides, we saw that especially \ac{TD3} and \ac{SAC} are unstable in the simultaneous self-play, where all agents are trained simultaneously (\autoref{sec:self-play}). Instead, the training against previously trained fixed checkpoints as self-play is more effective and used in the evaluation.

In conclusion, this report shows that the game factor of the hockey environment is more important than small algorithm modifications. The hockey environment tends to overfit to specific opponents, requiring numerous training runs and making performance heavily dependent on the choice of opponents. In the challenge, different opponents are used, such that a general policy is required to compete against the other teams. Thus, the biggest improvements are made through own rewards, fixed self-play, and long training runs, which help to train a general policy. Furthermore, our framework helps to explore more algorithms like CrossQ, hierarchical \ac{SAC}, or new hyperparameter combinations through mutations.



% Anpassungsfähig, overfitted aber schnell

% \begin{table}[h]
\caption{Performance between the agents and the baseline opponents.} \centering {}
\label{tab:performance_opponents}
\begin{tabular}{lccc}
\toprule
Agent & Hybrid Model & SAC (Own) & TD3 (Own) \\
Opponent &  &  &  \\
\midrule
\textbf{strong} & 94.60 & 97.10 & \textbf{97.90} \\
\textbf{weak} & 97.70 & 98.40 & \textbf{99.20} \\
\bottomrule
\end{tabular}
\end{table}

